\begin{statement}{C.2.UN.36}{Proved}{Compact component cores give a multiplicity-weighted invariant lattice in one nearby torus}
Let $f:U\to\Delta$ be a proper holomorphic map from
a smooth complex threefold, with connected two-torus
fibres over $\Delta^*$. Write its full central fibre
as $\sum_i m_iS_i$, and put $S=\bigcup_iS_i$ with
its reduced structure. For each component choose a
resolution of its normalization
$h_i:Z_i\to S_i\subset S$ which is an isomorphism
over $S_i\setminus\operatorname{Sing}S$. Set
\[
 \begin{gathered}
 D_i=h_i^{-1}(\operatorname{Sing}S)_{\rm red},\qquad
 R_i=\operatorname{span}_{\mathbb R}
       \{[C]:C\text{ is a curve component of }D_i\},
 \\
 K_i=\{v\in H_2(Z_i,\Z)/\operatorname{tors}:v\cdot R_i=0\}.
 \end{gathered}
                                                               \tag{C.CL.1}
\]
The set $D_i$ may also have isolated points; they
contribute no curve class to $R_i$. The pairing on
$K_i$ is the restriction of the ordinary oriented
surface intersection form and may have a radical.

Choose a common base-change order $N$ divisible by
all $m_i$ and by the semisimple monodromy order.
Let $U_1$ be the monodromy on the integral first
homology of the actual nearby torus after this
base change. There is a homomorphism, obtained by
choosing compact representatives and their transfers,
\[
 \Psi:\bigoplus_iK_i\longrightarrow
       H_2(F,\Z)^{\Lambda^2U_1},\qquad
 \Psi(v)\cdot\Psi(w)=\sum_i m_i(v_i\cdot w_i).
                                                               \tag{C.CL.2}
\]
Thus every multiplicity appears in the preserved
bilinear form. The kernel is contained in the
radical of that form. If its source is nondegenerate
of rank six, its image has finite integral index
in $H_2(F,\Z)$ and
\[
 \left|\det\left(\bigoplus_i m_iQ_{K_i}\right)\right|
                       \text{ is an integer square}.
                                                               \tag{C.CL.3}
\]

Let $(p_i,q_i,z_i)$ be the positive, negative and
radical dimensions of the real form on $K_i$.
Let $(p_U,q_U,z_U)$ be those of the intersection
form restricted to
$\ker(\Lambda^2U_1-I)\subset H_2(F,\mathbb R)$.
Then
\[
             \sum_i p_i\le p_U\le3,\qquad
             \sum_i q_i\le q_U\le3.                 \tag{C.CL.4}
\]
In particular, if $r_i=\dim R_i$, then
\[
 \sum_i\max\{b_2^-(Z_i)-r_i,0\}\le q_U\le3,
 \qquad
                 2\sum_i p_g(Z_i)\le p_U.           \tag{C.CL.5}
\]
A K3 component whose chosen resolution is obtained
from its minimal model by $n$ point blowups must
therefore have at least $16+n$ independent boundary
curve classes. For an Enriques component or a
properly elliptic component of coherent Euler
characteristic one, the corresponding lower bound
is $6+n$.

The existing monodromy theorem gives
$(U_1-I)^2=0$. In an oriented real basis of
$H_1(F,\mathbb R)$, all its possible Jordan forms
give the following exact invariant forms:
\[
\begin{array}{c|c|c|c}
\operatorname{rank}(U_1-I)&\text{nonzero values of }U_1-I
 &\dim\ker(\Lambda^2U_1-I)&(p_U,q_U,z_U)\\ \hline
0&\text{none}&6&(3,3,0)\\
1&e_2\mapsto e_1&4&(1,1,2)\\
2&e_2\mapsto e_1,\ e_4\mapsto e_3&4&(1,2,1)\\
2&e_2\mapsto e_1,\ e_4\mapsto-e_3&4&(2,1,1).
\end{array}                                                   \tag{C.CL.6}
\]
The two last rows record the two oriented real
forms of that Jordan type. In particular rank-one
unipotent monodromy forces $p_g(Z_i)=0$ for every
normalization component. All these assertions
apply to arbitrary actual central fibres, with
no normal-crossing, reducedness or K\"ahler
hypothesis on the original fibre.
\end{statement}

\begin{proof}
The open set
$Z_i^\circ=Z_i\setminus D_i$ maps isomorphically
to $S_i\setminus\operatorname{Sing}S$. It is
therefore precisely the part of this component
where the reduced support is one smooth embedded
surface. The compact analytic set $D_i$ consists
of curves and possibly isolated points. Its
degree-two integral cohomology is free, with
one generator for each curve component. Indeed
normalization of a compact reduced curve changes
only the identifications at finitely many points,
which do not change its degree-two cohomology;
isolated points have none.

Poincar\'e--Alexander duality and the relative
homology sequence give the exact map
\[
 H_2(Z_i^\circ,\Z)\longrightarrow H_2(Z_i,\Z)
 \longrightarrow H_2(Z_i,Z_i^\circ;\Z)
       \simeq H^2(D_i,\Z).                          \tag{C.CL.7}
\]
The second arrow, under these identifications,
is the vector of intersections with the curve
components of $D_i$. This follows by applying
Poincar\'e duality to the restriction map to
that compact analytic set. Torsion maps to
zero since the target is free. Consequently
every element of $K_i$ has an integral lift
represented by a compact finite cycle in
$Z_i^\circ$. Choose such lifts for an integral
basis of each $K_i$. Their mutual intersections
can be computed by compact transverse smooth
chain representatives in that same open set.
Include all the required perturbations in
one compact subset for each component.

Make the actual change $t=s^N$, normalize,
and resolve the total space and central support,
isomorphically over the smooth open sets just
specified and over $s\ne0$. Let
$g:\widetilde U\to\Delta_s$ be the resulting
proper map from a smooth threefold. Over a
point of $S_i\setminus\operatorname{Sing}S$
the original local equation is
$t=u x^{m_i}$, with $u$ a unit. Its normalized
root charts are
\[
       w^{m_i}=u,\qquad x=s^{N/m_i}/w.               \tag{C.CL.8}
\]
At $s=0$ the unit-root equation has exactly
$m_i$ local solutions, each nonsingular.
The full collection of strict components
above $S_i$ therefore restricts to a finite
\'etale covering
\[
        \pi_i:\widetilde Z_i^\circ
                       \longrightarrow Z_i^\circ
                  \quad\text{of total degree }m_i.  \tag{C.CL.9}
\]
It may be disconnected. Every one of these
strict components has multiplicity one;
the root normal coordinate is $s$.
No modification used in the resolution
changes (C.CL.8)--(C.CL.9).

Transfer an integral compact cycle by taking
the sum of all its lifts under (C.CL.9),
subdividing its finite singular chain into
evenly covered pieces. This is the ordinary
integral covering transfer and is a chain
map. For two transverse compact cycles,
every signed intersection point downstairs
has exactly $m_i$ lifts with the same local
orientation. Hence their transferred
intersection is exactly $m_i$ times the
original intersection. This proves the
multiplicity factor directly, including
disconnected root covers.

The transferred cycles lie in the smooth
single-component locus of the central
support of $\widetilde U$. The map $g$ is
a submersion on a neighbourhood of their
compact union. Choose smooth horizontal
lifts of the two real coordinate fields
on the base disc. A partition of unity
preserves their specified images under
$dg$. After a cutoff outside a slightly
larger neighbourhood, their successive
real and imaginary flows give, for every
sufficiently small $s$, an actual smooth
embedding of a neighbourhood of the compact
cycles into $g^{-1}(s)$. The image lies in
that exact fibre, and the embedding preserves
the complex orientation by continuity.
Choose disjoint ambient neighbourhoods for
the compact subsets in distinct strict
components, and the same small disc for
all flows. Their images then remain
disjoint in one nearby torus $F$.

Transport the selected transferred basis
cycles and extend by integer linearity.
This defines the homomorphism $\Psi$.
The transfer intersection calculation and
the disjoint neighbourhoods give the
bilinear equality of (C.CL.2). The choice
of lifts is part of this construction;
no canonical splitting of (C.CL.7) is
needed.

The image is invariant under the actual
nearby monodromy. In fact the two-flow
construction defines a smooth family of
each finite cycle over the \emph{whole}
small $s$-disc. Its restriction to a
circle in that disc returns to the
identical parametrized cycle. In any
local differentiable trivialization of
the proper smooth family over $s\ne0$,
this continuous family has a locally
constant integral homology class.
Thus the class returns unchanged under
Gauss--Manin transport. The actual
identification
$H_2(F,\Z)\simeq\Lambda^2H_1(F,\Z)$
intertwines transport, so its image is
in the group written in (C.CL.2).
Only these explicitly transported
cycles enter the assertion.

If $\Psi(v)=0$, the preserved pairing
with every $\Psi(w)$ shows that $v$
is in the radical of the source form.
Choose a real complement of that radical
on which the restricted form is
nondegenerate. Its image is injective
and has the same positive and negative
indices. Maximal positive and negative
subspaces of the target restricted form
therefore give (C.CL.4). If the integral
source is nondegenerate of rank six,
$\Psi$ is an injection of full-rank
free lattices into $H_2(F,\Z)$.
The latter has the unimodular form
$3H$, of determinant $-1$. If the
image has index $c$, the change-of-basis
determinant gives source determinant
$-c^2$. This proves (C.CL.3).

Intersect a maximal negative definite
subspace of $H_2(Z_i,\mathbb R)$ with
the codimension-$r_i$ subspace $K_i\otimes\mathbb R$.
Its dimension is at least
$b_2^-(Z_i)-r_i$, so its negative index
is at least the corresponding nonnegative
maximum in (C.CL.5). For the second
inequality there, the real and imaginary
parts of holomorphic two-forms give the
positive $2p_g(Z_i)$-dimensional subspace
proved in (C.PG.7)--(C.PG.8). Their
integrals on every boundary curve
vanish by type; their dual homology
classes therefore belong to
$K_i\otimes\mathbb R$ by (C.CL.7).
This proves (C.CL.5).

A K3 minimal surface has intersection
signature $(3,19)$; $n$ point blowups
add $n$ orthogonal classes of square
$-1$. For Enriques surfaces the
signature is $(1,9)$. A properly
elliptic component with $\chi(\OO)=1$
has $q=p_g=0$ and $e=12$ by UN33,
and signature $(1,9)$ as well:
$b_1=0$, $b_2=10$, and the signature
formula with $K^2=0$ gives signature
$-8$. Applying the negative-index
bound to $19+n$ or $9+n$ gives
the asserted boundary ranks.

It remains to compute the exact
monodromy bound. C.2.UN.11 proves
that $N_0=U_1-I$ has square zero
and rank at most two. In rank one,
choose a vector and its nonzero
image, then complete the basis in
the kernel. Changing the sign of
one unused basis vector makes the
basis positively oriented, giving
the second row of (C.CL.6).
In rank two, choose two image
vectors and their preimages. If
their resulting ordered basis is
negative, reverse one preimage.
This gives the two oriented forms
\[
          N_0e_2=e_1,\qquad N_0e_4=\epsilon e_3,
                     \qquad\epsilon\in\{1,-1\}.     \tag{C.CL.10}
\]
Write $e_{ij}=e_i\wedge e_j$.
In the oriented wedge basis
$(e_{12},e_{13},e_{14},e_{23},e_{24},e_{34})$
the only nonzero complementary
pairings are
\[
 e_{12}\cdot e_{34}=1,\qquad
 e_{13}\cdot e_{24}=-1,\qquad
 e_{14}\cdot e_{23}=1.                              \tag{C.CL.11}
\]
For rank one the invariant space
has basis $(e_{12},e_{13},e_{14},e_{34})$;
its form is one hyperbolic plane
and a two-dimensional radical.
For (C.CL.10) an invariant basis is
\[
       e_{12},\ e_{34},\ e_{13},\ e_{14}-\epsilon e_{23}.
                                                               \tag{C.CL.12}
\]
Its Gram matrix is
\[
 \begin{pmatrix}
 0&1&0&0\\
 1&0&0&0\\
 0&0&0&0\\
 0&0&0&-2\epsilon
 \end{pmatrix}.                                               \tag{C.CL.13}
\]
These are exactly the rows of
(C.CL.6); the identity row is
the whole form $3H$.
The exterior-square matrices,
their exact kernels, and all
four invariant inertias are
executed over the integers and
rationals in
\path{code/A_compact_core_lattices.py}
and its saved output. Finally
the rank-one row has $p_U=1$,
so (C.CL.5) forces every
$p_g(Z_i)$ to vanish, as claimed.
\end{proof}


% C.2.D4.1 from part_C2_cubic_cone.tex

